
\section{Specifica dei Requisiti del Sistema}

\subsection{Requisiti di Business (BR)}

\begin{description}[leftmargin=0.5cm]
    \item[BR1 (Scopo e Successo):] Il sistema fornirà un ambiente desktop di live coding musicale. Il progetto sarà valutato come un successo se permetterà a un utente di scrivere istruzioni ritmiche e melodiche,
    riprodurre audio in tempo reale, il tutto con una stabilità.
    \item[BR2 (Deadline):] Implementare il suddetto applicativo in un tempo limite di due mesi.
\end{description}

\subsection{Requisiti Funzionali:}
\subsubsection{Requisiti Utente (RU)}

\begin{description}[leftmargin=0.5cm, itemsep=0.8em]
    \item[RU1 (Editor di Codice):] Il sistema deve fornire un'area di testo multiriga (Editor) in cui l'utente può inserire, modificare e formattare liberamente sequenze di testo.

    \item[RU2 (Gestione Esecuzione):] L'utente deve disporre di un comando esplicito (pulsante visibile nell'interfaccia o scorciatoia da tastiera) per gestire il codice attualmente scritto nell'Editor.
        \subitem Button play (ctrl+enter)
        \subitem Button stop (ctrl+.)
        \subitem Button update (ctrl+u)


\item[RU11 (Sintassi base):] L'utente ha a disposizione due metodi di base che sono:
        \begin{itemize}
            \item \texttt{note()}: accetta come parametro una stringa con all'interno le note che dovrà riprodurre. Le note possono essere corredate da modificatori (specificati più avanti).
            \item \texttt{sound()}: serve a riprodurre file audio campionati o a impostare il tipo di strumento. Di default è impostato su suoni di batteria (Figure~\ref{fig:etichetta0}).
        \end{itemize}

    \item[RU12 (Sequenziamento Base):]
    L'utente genera una sequenza ritmica inserendo una stringa testuale composta da nomi di campioni audio o note separati da spazi vuoti (es. \textbf{sound("bd hh sd hh")} oppure \textbf{note("48 52 55 59").sound("piano")}).
    Il sistema non usa una griglia a step fissi, ma costringe l'intera sequenza all'interno di un "ciclo" di tempo predefinito (per default la durata di un ciclo è di due secondi), dividendo matematicamente e in parti uguali la durata del ciclo per il numero di suoni inseriti.
    Utilizzando le parentesi angolari (es. \textbf{sound("<bd hh rim>")}), il sistema eseguirà \textbf{un solo elemento} di questo blocco per ogni ciclo, alternandoli ciclicamente a ogni nuova ripetizione.
    NB. se usiamo una sequenza di questo tipo "bd <hh oh hh bs>*2" il sistema dividerà il ciclo in due metà: nella prima metà suonerà in modo fisso il campione \textbf{bd}, mentre nella seconda metà, grazie al moltiplicatore *\textbf{2}, il blocco di alternanza verrà valutato due volte (avanzando di due step alla volta). Di conseguenza, nel primo ciclo l'output temporale effettivo sarà equivalente alla sequenza \textbf{bd [hh oh]} (dove \textit{hh} e \textit{oh} si dividono lo spazio della seconda metà del ciclo), nel secondo ciclo sarà equivalente a \textbf{bd [hh bs]}, e dal terzo ciclo il pattern ricomincerà da capo


    \begin{figure}[htbp]
        \centering
        \includegraphics[width=0.8\textwidth]{drumSound}
        \caption{Drum sound}
        \label{fig:etichetta0}
    \end{figure}

    \item[RU13 (Variazioni e Pause):]
    L'utente può creare "silenzi" o sincopi all'interno della sequenza inserendo i caratteri speciali \texttt{-} (trattino) o \texttt{\textasciitilde} (tilde) separati da spazi.
    Il sistema conta questi caratteri come elementi validi nel calcolo della suddivisione del ciclo: la pausa occupa l'esatta frazione di tempo di un suono normale, ma senza innescare alcuna riproduzione audio.


    \item[RU13 (Raggruppamenti e Moltiplicatori):] L'utente deve poter raggruppare eventi in sotto-sequenze tramite le parentesi quadre (es. \texttt{bd [hh hh] sd}), costringendoli a dividersi lo spazio di un singolo step. Deve poter alterare la velocità di un singolo step o di un gruppo utilizzando l'asterisco (es. \texttt{*2} per raddoppiare la velocità, \texttt{/2} per dimezzarla).
    \item[RU14 (Parallelismo):] L'utente deve poter eseguire più pattern in perfetta sovrapposizione simultanea separandoli con una virgola all'interno della stringa (es. \texttt{bd bd, hh hh hh hh}).



    \item[RU16 (Notazione Anglosassone e Ottave):] L'utente deve poter comporre melodie scrivendo direttamente le note tramite la notazione alfabetica standard (es. \texttt{note("c d e f")}). Deve poter specificare l'altezza aggiungendo il numero dell'ottava (es. \texttt{c4}) e alterare l'intonazione usando \texttt{\#} (diesis) o \texttt{b} (bemolle).


    \item[RU18 (Catena di Effetti Audio):] L'utente deve poter alterare il timbro del pattern concatenando metodi con la notazione a punto. Deve avere a disposizione comandi per il posizionamento stereo (\texttt{.pan()}), il volume (\texttt{.gain()}), riverberi spaziali (\texttt{.room()}), echi (\texttt{.delay()}), e filtri di frequenza (\texttt{.lpf()}, \texttt{.hpf()}).


    \item[RU19-OPZIONALE (Modulazione Dinamica dei Parametri):] L'utente non deve essere limitato a inserire valori numerici statici negli effetti (es. \texttt{.pan(0.5)}), ma deve poter inserire intere stringhe in Mini Notation come parametri (es. \texttt{.pan("<0 1>")}). Questo permette all'utente di far muovere o evolvere l'effetto nel tempo in modo sincronizzato al ritmo.



    \item[RU20-OPZIONALE (Inversione e Frammentazione):] L'utente deve poter manipolare l'ordine e la densità degli eventi generati usando specifiche funzioni. Deve poter suonare un pattern al contrario (es. usando \texttt{.rev()}) o suddividere un singolo evento in microscopiche ripetizioni (es. usando \texttt{.ply(3)}).
    \item[RU21-OPZIONALE (Spazializzazione e Sfasamento):] L'utente deve poter creare poliritmie complesse sdoppiando un pattern. Deve poter usare funzioni come \texttt{jux()} per mandare una versione processata del pattern sul canale audio destro e l'originale sul sinistro, oppure usare \texttt{off()} per sovrapporre al pattern una copia di se stesso ritardata nel tempo.


    \item[RU22-OPZIONALE (Sintesi e Generazione Continua):] L'utente deve poter sostituire la riproduzione di campioni audio discreti con sintetizzatori generativi (es. onde sinusoidali o a dente di sega tramite comandi come \texttt{s("sawtooth")}), permettendo la creazione di suoni continui modellati da inviluppi matematici.



    \item[RU23 (Dilatazione e Compressione):] L'utente deve poter alterare la velocità globale di interi blocchi di codice applicando modificatori come \texttt{.fast(2)} (per raddoppiare la frequenza di esecuzione) o \texttt{.slow(2)} (per dimezzarla).
    \item[RU24 (Anticipo e Ritardo):] L'utente deve poter spostare l'intero pattern in avanti o indietro rispetto al metronomo principale, utilizzando funzioni di sfasamento (es. \texttt{.early()} o \texttt{.late()}).


    \item[RU4-OPZIONALE (Feedback Visivo di Riproduzione)]
    Durante l'esecuzione dell'audio, l'utente deve vedere la specifica parola/token all'interno dell'Editor di testo evidenziata (es. tramite cambio di colore dello sfondo del testo) esattamente nell'istante in cui il suono corrispondente viene emesso dalle casse.
    \begin{figure}[htbp]
        \centering
        \includegraphics[width=0.8\textwidth]{patternColorati}
        \caption{Esempio di visualizzazione dei pattern in esecuzione}
        \label{fig:etichetta}
    \end{figure}


    \item \textbf{RS5-OPZIONALE (Area di Visualizzazione):}
    \subitem \textbf{Pianoroll}
    \newline
    L'interfaccia utente deve includere un'area grafica bidimensionale,
    posizionata adiacente o sottostante all'Editor di codice, dedicata alla rappresentazione
    del "Piano Roll" (o Punchcard).


    Il sistema deve mappare gli eventi sonori sull'area grafica secondo due assi:

    \subsubitem \textbf{L'Asse X (Orizzontale) rappresenta il Tempo:}
     Lo spazio orizzontale disponibile deve essere diviso proporzionalmente
     in base alle frazioni di ciclo calcolate dal parser. L'inizio
     e la fine di un evento sull'asse X devono corrispondere esattamente
     al suo inizio e fine temporale all'interno del pattern.
    \subsubitem \textbf{L'Asse Y (Verticale) rappresenta l'Altezza (Pitch/Nota):}
     Il sistema deve mappare il valore della nota musicale (es. c, a, f, e)
     su un'altezza verticale specifica. Note più acute (frequenze maggiori)
     devono essere posizionate più in alto sull'asse Y; note più gravi
     (frequenze minori) più in basso.


    Ogni evento sonoro deve essere disegnato come un blocco rettangolare continuo all'interno
    dell'area grafica. La larghezza del blocco deve essere proporzionale alla durata della nota.


    Il sistema deve renderizzare un indicatore verticale (es. una linea continua)
    che si sposta orizzontalmente attraverso l'area grafica da sinistra
    verso destra, rappresentando il tempo di riproduzione corrente.


    Il movimento della testina di lettura  deve essere fluidamente animato
    e mantenuto in stretta sincronia  con il motore di esecuzione audio.


    Quando la testina di lettura interseca il perimetro orizzontale
    di un blocco rettangolare (evento sonoro), il sistema deve alterare lo
    stato visivo di quel blocco (es. evidenziandone i bordi o cambiandone colore)
    per indicare chiaramente l'istante in cui la nota sta suonando attivamente,
    ripristinando lo stato precedente quando la testina supera il blocco.

    \begin{figure}[htbp]
        \centering
        \includegraphics[width=0.8\textwidth]{pianoroll}
        \caption{Esempio di visualizzazione de`l piano roll`}
        \label{fig:etichetta1}
    \end{figure}

    \subitem \textbf{Oscilloscopio}
    \newline


    Il sistema deve fornire un componente visivo, richiamabile dall'utente tramite i comandi \_scope() o tscope(),
    dedicato a renderizzare un oscilloscopio che mostri il segnale audio nel dominio del tempo. Tale visualizzazione
    deve apparire nell'interfaccia adiacente o sottostante all'editor di codice.


    \textbf{(Rappresentazione della Forma d'Onda):} Il sistema deve mappare i campioni audio su un grafico
    bidimensionale disegnando una linea continua. In questo grafico, l'asse X orizzontale rappresenta lo
    scorrere del tempo, mentre l'asse Y verticale rappresenta il valore dell'ampiezza istantanea del segnale
    audio (es. una forma d'onda a dente di sega come sawtooth).


    \textbf{(Parametri di Stile Grafico):} Il parser e il motore di rendering devono
    supportare un oggetto di configurazione per permettere all'utente di alterare dinamicamente
    l'aspetto della linea tracciata, supportando i seguenti attributi:

    \begin{itemize}
        \item\textbf{Colore (color):} Accetta stringhe rappresentanti colori (esadecimali o nomi standard).
    Il sistema deve applicare il bianco (white) come valore predefinito.


    \item \textbf{Spessore (thickness):} Accetta un valore numerico che definisce lo spessore della linea in pixel (o unità equivalenti). Il sistema deve applicare un valore predefinito pari a 3.


    \item \textbf{Scala (scale):} Accetta un valore numerico che agisce come moltiplicatore di zoom per l'asse Y, ingrandendo o riducendo verticalmente la forma d'onda. Il sistema deve applicare un valore predefinito pari a 0.25.


    \item \textbf{Posizione (pos):} Accetta un valore numerico per traslare l'intera forma d'onda lungo l'asse verticale dell'area di visualizzazione, dove il valore 0 allinea l'onda al margine superiore dello schermo e il valore 1 la allinea al margine inferiore.


    \end{itemize}

    \textbf{Algoritmo di Stabilizzazione e Trigger:} Per garantire che l'onda visualizzata risulti stabile
    all'occhio umano (evitando l'effetto di scorrimento casuale tra i frame di aggiornamento), il motore
    grafico deve implementare un meccanismo di triggering basato sui seguenti parametri:

    \begin{itemize}
    \item \textbf{Allineamento (align):} Un parametro booleano. Se impostato a 1 (vero), il sistema deve iniziare
    a disegnare il frame visivo partendo esattamente dal punto in cui il segnale interseca il valore di trigger
    (zero-crossing). Questo comportamento deve essere attivo (1) di default.


    \item\textbf{Ampiezza di Trigger (trigger):} Un valore numerico di ampiezza che il sistema deve utilizzare
    come soglia per attivare la partenza del disegno della linea. Il valore predefinito deve essere 0.


    \end{itemize}

    \begin{figure}[htbp]
        \centering
        \includegraphics[width=0.8\textwidth]{scope}
        \caption{Esempio di visualizzazione dello scope}
        \label{fig:etichetta2}
    \end{figure}

\item \end{description}

\subsubsection{Requisiti di Sistema (RS)}

\vspace{0.5cm}
\noindent\textbf{Modulo Parser e Interprete}
\begin{description}[leftmargin=0.5cm, itemsep=0.8em]
    \item[RS1 (Parsing della Stringa):] Alla ricezione del comando di Avvio (RU2), il sistema deve scansionare la stringa dell'Editor e generare una struttura dati gerarchica (es. \emph{Abstract Syntax Tree}) che rappresenti l'ordine, il raggruppamento (es. [ ]), i moltiplicatori (es. *2) e le sequenze parallele (es. ,). L'algoritmo di parsing deve rifiutare stringhe malformate restituendo un messaggio di errore chiaro senza far crashare il sistema.

    \item[RS2 (Calcolo delle Frazioni di Tempo):] Il parser deve calcolare con precisione matematica la frazione di ``ciclo base'' occupata da ogni singolo token testuale, suddividendo equamente lo spazio tra gli elementi di una sequenza o di un sottogruppo.

    \item[RS3 (Mappatura Posizionale):] Durante il parsing, il sistema deve salvare e associare a ogni evento sonoro calcolato le coordinate esatte (indice di inizio e fine carattere) della parola originale corrispondente all'interno dell'Editor di testo.
\end{description}

\vspace{0.5cm}
\noindent\textbf{Motore Audio e Sincronizzazione}
\begin{description}[leftmargin=0.5cm, itemsep=0.8em]
    \item[RS4 (Scheduling degli Eventi):] Il sistema deve tradurre le frazioni temporali relative in millisecondi assoluti (basandosi su un parametro \emph{Cycles Per Minute} configurabile) e programmare l'esecuzione dei campioni audio nel futuro per evitare ritardi di elaborazione.

    \item[RS5 (Notifica di Avvenuta Esecuzione):] Nel momento esatto in cui il motore audio innesca la riproduzione di un campione, il sistema deve emettere un segnale di notifica contenente le ``coordinate testuali'' calcolate nel requisito RS3.

    \item[RS6 (Aggiornamento della UI):] L'interfaccia grafica deve intercettare le notifiche descritte in RS5 e applicare dinamicamente lo stile di ``evidenziazione'' esclusivamente ai caratteri compresi in quelle coordinate, ripristinando lo stile normale non appena l'evento sonoro ha terminato la sua durata teorica.
\end{description}

\subsection{Requisiti Non Funzionali (RNF)}
Stabiliscono le qualità del sistema e devono essere oggettivamente validabili tramite test.

\begin{description}[leftmargin=0.5cm, itemsep=0.8em]
    \item[RNF1 (Architettura Multi-Thread):] Per garantire la reattività dell'interfaccia, il processo di calcolo audio (DSP e Scheduling) e il processo di aggiornamento dell'interfaccia grafica (GUI) devono necessariamente operare su thread separati. Il congelamento (\emph{freeze}) del thread grafico non deve mai interrompere il flusso audio.

    \item[RNF2 (Latenza Visiva):] Il ritardo tra l'effettiva emissione dell'audio (misurata internamente al motore sonoro) e l'aggiornamento grafico del token nell'Editor (RS6) non deve superare i 30 millisecondi, garantendo un feedback percepito come ``istantaneo'' dall'occhio umano.

    \item[RNF3 (Jitter Audio):] Il sistema di scheduling (RS4) deve garantire una deviazione massima (\emph{jitter}) dal tempo ideale di esecuzione di un suono inferiore ai 5 millisecondi, per prevenire aritmie udibili.

    \item[RNF4 (Tolleranza ai Cambiamenti):] La modifica del testo nell'Editor durante l'esecuzione dell'audio (live coding vero e proprio) non deve bloccare il suono in corso, ma deve essere programmata per avere effetto a partire dal successivo ``ciclo'' musicale.
\end{description}

\subsection{Requisiti di Implementazione (RI)}
Vincoli specifici per la fase di produzione.

\begin{description}[leftmargin=0.5cm, itemsep=0.8em]
    \item[RI1 (Tecnologia di Sviluppo):] Il sistema sarà sviluppato interamente in Scala.

    \item[RI2 (Isolamento del Dominio):] Il motore di parsing (RS1, RS2) deve essere scritto come modulo software puro, totalmente slegato da librerie di interfaccia grafica o da librerie audio specifiche, affinché sia testabile in totale isolamento.
\end{description}
